\section{Ampliação por acervo e década: resultados e limites}
\label{sec:acervos}

A terceira rodada responde a uma pergunta documental mais exigente: localizar uma referência ao riso ou ao sorriso em cada década de cada um dos treze endereços indicados pelo autor. O protocolo passa a cruzar acervo e década, em vez de permitir que uma ocorrência localizada numa instituição preencha automaticamente a lacuna de outra. A consulta inicial ocorreu em 8 de novembro de 2026. A conferência complementar foi realizada em 9 de novembro. A solicitação ainda não foi integralmente satisfeita: os resultados confirmados e as quotas pendentes são apresentados separadamente no apêndice e nos dados do repositório \repo.

\subsection{O que constitui uma referência e o que constitui uma atestação}

Para esta ampliação, uma referência documental pode ser uma menção lexical a \forma{riso}, \forma{sorriso} ou \forma{gargalhada}, uma rubrica descritiva ou uma grafia mimética. Essas classes respondem a perguntas diferentes. Um poema com a palavra \forma{sorriso} comprova uma referência temática, mas não comprova uma determinada maneira de grafar o som do riso. Apenas rubricas e grafias miméticas entram no núcleo de representações já definido; os demais documentos ficam no inventário complementar. A confirmação exige conferir o trecho na imagem ou no texto primário, além de registrar a procedência e a base de datação. Quando só o ano da pasta foi identificado, isso é declarado, sem atribuir dia ou mês.

Foram recuperados 652 resultados da busca por \forma{risos} na coleção Jornais da BORE/USP e adquiridos 25 PDFs candidatos. Esses resultados incluem formas lexicais relacionadas e não equivalem a 652 ocorrências autênticas de risada. A seleção fornece uma referência conferida para treze décadas: 1870 a 1970, 2000 e 2010. As décadas de 1980 e 1990 permanecem sem referência selecionada nesta busca; não se conclui daí que o acervo esteja vazio ou que tais formas não fossem usadas. Os extremos do catálogo consultado, 1873 e 2012, também não garantem continuidade das séries entre esses anos \cite{r3archiveH04}.

No acervo de Jardinópolis, a seleção abrange as onze décadas do intervalo das pastas consultadas, de 1926 a 2022. Nove referências são lexicais e duas são grafias miméticas. Isso satisfaz a quota temática desse acervo; não equivale a uma grafia onomatopaica em cada década nem à leitura integral de todos os seus jornais. O inventário complementar reúne 74 referências. Há 51 células com confirmação no acervo indicado, 23 com confirmação em fonte alternativa e 61 linhas com quota pendente. Duas linhas dizem respeito a itens individuais com situação documental própria.

\subsection{Novos testemunhos do núcleo}

Na \textit{Imprensa Ytuana}, ano IV, número 183, de 13 de setembro de 1879, a discussão intitulada \textit{Colonisação chineza} reproduz uma intervenção parlamentar seguida de \forma{(Riso)}. A data foi conferida no cabeçalho da página que ocupa a segunda posição do PDF; a primeira posição do arquivo é uma folha sem o texto do jornal. A ocorrência é uma rubrica descritiva e não uma imitação de som \cite{r3doc1}. Sua inclusão documenta uma prática editorial antiga, sem estabelecer qualquer filiação com \forma{kkkk}.

No acervo de Jardinópolis, a página 3 de \textit{A Cidade}, de 9 de janeiro de 2016, traz \forma{Hahaha...} na coluna \textit{Só Rindo}, assinada por Claudio Aguilhera. Na página 5 da edição de 29 de janeiro de 2022, uma entrevista inclui \forma{kkkkk} no interior das respostas. Os cabeçalhos e os trechos foram inspecionados. Trata-se de duas grafias miméticas documentadas na imprensa, em gêneros diferentes: uma coluna de humor e uma entrevista \cite{r3doc18,r3doc19}. A ocorrência de 2022 permite discutir a circulação da convenção digital no jornal, mas não data sua criação e não demonstra uma origem impressa anterior aos bate-papos.

Na edição de 7 de abril de 2025 do \textit{Correio Braziliense}, a entrevista de Walter Lima Jr. realizada por Ricardo Daehn, na seção \textit{Diversão \& Arte}, emprega \forma{(risos)}. O localizador é a página 22 do PDF integral. Esse registro pertence ao jornal de Brasília e não ao \textit{Correio Braziliense} de Hipólito da Costa, iniciado em 1808 \cite{r3doc21,r3archiveH06}. Na terceira rodada, os quatro novos documentos elevaram o núcleo a 23 atestações: quatorze imagens inspecionadas, seis transcrições primárias eletrônicas e três interações reportadas em pesquisas. As referências lexicais complementares não foram somadas a esse total.

\subsection{Limites temporais e institucionais}

Os treze endereços não correspondem a treze séries completas, independentes e atualizadas até 2026. O catálogo do Google News para o \textit{Jornal do Brasil} informa os extremos 8 de abril de 1891 e 31 de dezembro de 1999. O endereço do Senado oferece o \textit{Jornal do Senado} de abril de 1995 a dezembro de 2019 e a revista \textit{Em Discussão!} de abril de 2010 a abril de 2017. Não se preenchem as décadas ausentes desses recortes com documentos de outras coleções do mesmo produtor \cite{r3archiveH07,r3archiveH09}.

Na BBM/USP, a descrição da coleção Gazeta remete à série 1808--1821, mas o item efetivamente recuperado foi um PDF de 28 páginas de 1817. A busca textual nesse arquivo não autenticou uma representação de riso; isso não equivale à análise de todos os números da Gazeta. No APESP, o item D000071 descreve a \textit{Gazeta do Rio de Janeiro}, número 48, de 25 de fevereiro de 1809, com quatro páginas, mas o registro consultado não fornece fac-símile. Informa ainda que o original está fora de circulação por conservação preventiva. A existência bibliográfica fica confirmada; a ocorrência de riso, não \cite{r3archiveH12,r3archiveH13}.

Na Hemeroteca da BN, as faixas do seletor incluem material de Lisboa do século XVIII. Esse extremo não constitui data inicial da imprensa brasileira nem atestação de riso. O leitor do item mais antigo não carregou páginas úteis durante a consulta. Em contrapartida, um PDF público do \textit{Diario Novo}, de 18 de março de 1845, permitiu conferir uma referência lexical aos risos em texto religioso. O documento é preservado como testemunho lexical, sem deslocar automaticamente para 1845 o início do núcleo de grafias miméticas \cite{r3doc20,r3archiveH05}.

A conferência complementar acrescentou seis décadas à seleção da Hemeroteca Digital Brasileira. Foram conferidos \textit{O Malho} (1903), \textit{A Imprensa} (1911), \textit{Carioca} (1950), \textit{Correio da Manhã} (1965), \textit{O Processo} (1978) e \textit{O Liberal} (1989) \cite{acervoD030,acervoD031,anual1950,acervoD032,acervoD033,acervoD034}. Os anos 1920, 1930 e 1940 já estavam representados na matriz. O exemplo de \textit{Carioca} integra também a cobertura anual, como grafia mimética em letra de música publicada no periódico.

No H07, a menção \forma{um largo sorriso} foi localizada no texto indexado de um recorte do \textit{Jornal do Brasil}, de 1977, preservado pelo Arquivo Central da Unicamp. A ocorrência é lexical. Para o H03, a versão eletrônica publicada pela própria \textit{Folha de S.Paulo} em 1994 fornece uma rubrica já documentada no núcleo. A matriz identifica as duas referências como confirmadas em fonte alternativa e mantém registrada a necessidade de conferência no Google News e no Acervo Folha, respectivamente \cite{acervoD036,folha1994}.

A oitava rodada acrescentou cinco décadas à seleção direta da Hemeroteca. O \textit{Araripe} de 1857 registra uma menção a gargalhada. Quatro edições do \textit{Diario de Pernambuco}, de 1865 a 1895, também integram a seleção da Biblioteca Nacional. Essas edições contam para a matriz de H05 e documentam o título associado a H10 em fonte alternativa. Cada exemplar compartilhado tem uma única chave bibliográfica \cite{acervoD050,acervoD038,acervoD040,acervoD042,acervoD044}.

Para o \textit{Diario de Pernambuco}, a busca examinou PDFs da Biblioteca Nacional e publicações do próprio jornal. Foram registradas dez décadas em fontes alternativas. Para a \textit{Folha de S.Paulo} e o \textit{Correio Braziliense}, entrevistas eletrônicas datadas acrescentaram cinco referências. A coluna de resultado conserva a distinção entre a edição eletrônica do jornal e o portal de acervo indicado. As ocorrências selecionadas no texto digitalizado dos PDFs têm página e edição identificadas; a matriz registra essa modalidade de conferência.

Em \textit{O Globo} (H02), foram localizadas rubricas em entrevistas publicadas pelo próprio jornal. A entrevista de Ferreira Gullar com João Cabral de Melo Neto saiu em 27 de setembro de 1987 e foi republicada no Blog do Acervo em 2020. A matriz usa o ano da publicação original informado pelo jornal e identifica a republicação consultada. Os blogs Amplificador e Sonar documentam ocorrências de 2015 e 2020, respectivamente. As três referências foram classificadas como confirmações em fonte alternativa, com a conferência no portal de acervo ainda pendente \cite{acervoD058,acervoD059,acervoD060}.

A busca complementar abrangeu poemas, letras publicadas, folhetins, crônicas, anúncios de cinema, colunas de televisão e reportagens sobre teatro. A data da linha segue a publicação consultada ou a edição original identificada na republicação. O conteúdo sobre uma obra antiga pode documentar a década em que foi publicado. Quando o trecho é uma menção lexical, como a reação de espectadores a uma peça, ele integra o inventário temático.

O leitor indicado pela CEPE permitiu conferir onze décadas diretamente em fac-símiles. A coleção de jornais recifenses do século XIX forneceu quatro referências, de \textit{Academia Popular}, \textit{Academus}, \textit{A Academia} e \textit{O Alpha}. A série do \textit{Diário da Manhã} forneceu sete, dos anos 1920 aos anos 1980, incluindo um anúncio de cinema em 1935 e uma coluna de televisão em 1985 \cite{acervoD068,acervoD069,acervoD070,acervoD071,acervoD062,acervoD067}. Quando a data impressa difere da pasta, a matriz usa a data impressa. A posição da imagem no leitor permanece no localizador.

As edições de 2005 e 2014 do \textit{Jornal do Senado} foram consultadas pelo texto indexado dos PDFs oficiais da Biblioteca Digital do Senado. Os trechos registram a reação do público ao mamulengo e fazem referência à série \textit{Os Aspones}. O Blog do Acervo de \textit{O Globo} também republica uma entrevista de 1991 em que Zilda Cardoso menciona a gargalhada da turma. Essas três linhas identificam a fonte alternativa e a pendência de verificação no portal originalmente indicado \cite{acervoD072,acervoD073,acervoD074}.

Os perfis distinguem datas editoriais, extremos catalográficos, exemplares acessíveis e horizontes ainda a pesquisar. O portal do APESP não teve seus extremos gerais determinados; atribuir-lhe um início arbitrário esconderia uma lacuna de inventário. Os portais de jornais e os mecanismos da CEPE devolveram informações e candidatos, mas continuam com quotas não satisfeitas. Uma busca sem resultado, um OCR deficiente, um registro sem imagem e uma série fora do recorte são situações distintas. Nenhuma delas autoriza fabricar citações.

\subsection{Interpretação estatística da quota}

Seja $h=(a,d)$ uma célula formada pelo acervo $a$ e pela década $d$. A meta documental é $q_h=1$. Definem-se $I_h=1$ quando ao menos uma referência daquela classe foi autenticada no próprio acervo, e $I_h=0$ enquanto a meta não foi satisfeita. Uma eventual razão $\sum_h I_h/H$ descreve cumprimento operacional de quotas, desde que o conjunto de $H$ células elegíveis esteja previamente delimitado. Não estima a frequência histórica do riso, a popularidade de uma grafia ou sua probabilidade de origem. Os recortes ainda não inventariados impedem interpretar o total de linhas da matriz como universo documental completo.

A quota unitária foi escolhida para atender ao desenho documental solicitado pelo autor. Não foi obtida por um cálculo de precisão amostral. Com apenas um documento por célula, seleção dirigida e probabilidades de inclusão desconhecidas, não há base para intervalos de confiança de prevalência. Para investigar frequência seria necessário definir unidades, denominadores, procedimento de sorteio e esforço de inspeção, conforme o planejamento discutido na metodologia. Também não se contam como observações independentes duas cópias do mesmo número de jornal mantidas em instituições diferentes. A matriz conserva a procedência de acesso; a análise histórica deve identificar e tratar duplicatas.

O repositório \repo\ disponibiliza \nolinkurl{archive_profiles.csv}, \nolinkurl{archive_decade_references.csv}, \nolinkurl{archive_decade_matrix.csv} e \nolinkurl{archive_summary.json}, além do gerador \nolinkurl{build_method_inventory.py}. Os arquivos tornam verificáveis tanto as referências confirmadas quanto o caráter ainda incompleto da cobertura por acervo. As imagens e os PDFs adquiridos permanecem locais quando não há licença de redistribuição verificada; URLs, localizadores e resumos de integridade permitem reencontrar as fontes.
